\scanpage{185}
\mspage{496}
% position: top continuation
le déséquilibre marqué par la valorisation exclusive du yang, auquel il m'est
arrivé de faire allusion déjà précédemment (*). La même constatation peut se
faire pour la quasi-totalité des couples yin-yang sur lesquels je suis tombé -
c'est là une chose vraiment très frappante, que je n'avais jamais pris le loisir
précédemment de vérifier de façon aussi circonstanciée.

% position: upper paragraph
Parmi les couples écrits tantôt, le seul qui me semble faire exception
est le couple \mbox{p a s s i o n - s é r é n i t é}, vu que dans l'usage courant\inserted{,}
le mot ``passion'' souvent s'associe à l'image d'un déchaînement, d'une violence,
ou sinon d'un \mbox{l a i s s e r - a l l e r}, avoisinant fâcheusement le nuage
d'associations entourant un mot comme ``\mbox{t u r p i t u d e}''. Comme par hasard,
laisser-aller et turpitude désignent des états de déséquilibre psychique carac-
térisés par une excessive prépondérance \mbox{y i n}, féminine! Et symétriquement,
suivant les mêmes mécanismes pousse-bouton (qui révèlent nos conditionnements
courants, et nullement la nature d'une chose comme la ``sérénité''), le mot ``sé-
rénité'' s'associe (par opposition à ``passion'') à l'image d'un \mbox{c o n t r ô l e}
\mbox{d e\quad s o i} - d'une qualité donc qui, comme de juste, est d'essence \mbox{m a s c u-}
\mbox{l i n e}. (En fait, le pendant yin du ``contrôle'' n'est nullement la ``passion'',
mais ``l'abandon''.)

% position: middle paragraph and displayed pair
Ce qui se passe donc ici, c'est que par suite d'une confusion générale
dans les esprits sur la nature de certaines choses, s'exprimant par une égale
confusion dans l'usage de certains mots, censés les désigner, il y a une con-
fusion du couple yang-yin ``passion-sérénité'' avec l'ensemble des deux notions
\begin{center}relâchement - contrôle,\end{center}
dont les termes sont yin-yang (sans pour autant constituer un ``couple'', les
deux termes en présence n'ayant aucune envie de s'épouser!). Il me semble donc
que la soi-disante ``exception'' à la règle (de la valorisation systématique du
yang) en est au contraire une confirmation particulièrement intéressante! Et
je ne serais pas étonné qu'il en soit de même avec les autres quelques exemples
que j'ai relevés, où dans un couple yang-yin, c'est le terme yin qui semblerait
valorisé.

% position: lower paragraph
Je ne suis d'ailleurs nullement sûr que cette distorsion dans la vision
du monde que je constate dans la civilisation dite ``occidentale'', provenant de ce
parti-pris systématique en faveur du masculin, opposé au féminin - que cette
distorsion, ce déséquilibre soient tellement moindres dans la tradition chinoise,
% position: footnote
\par\noindent\rule{0.24\linewidth}{0.3pt}\par
(*) Voir la note ``Yang enterre yin (1) - ou le muscle et la tripe'', n° 106.
\margin{184}
